\documentclass[10pt]{article}
\ifdefined\pdfinfoomitdate\pdfinfoomitdate=1\fi
\ifdefined\pdftrailerid\pdftrailerid{}\fi
\ifdefined\pdfsuppressptexinfo\pdfsuppressptexinfo=15\fi
\usepackage[margin=0.72in]{geometry}
\usepackage{amsmath,amssymb,amsthm,mathtools,booktabs,microtype,hyperref}
\usepackage[T1]{fontenc}
\setlength{\parskip}{2.5pt}
\setlength{\parindent}{1.2em}
\setlength{\abovedisplayskip}{6pt plus 2pt minus 2pt}
\setlength{\belowdisplayskip}{6pt plus 2pt minus 2pt}
\raggedbottom
\newtheorem{theorem}{Theorem}
\newtheorem{lemma}{Lemma}
\newtheorem{corollary}{Corollary}
\newcommand{\Z}{\mathbb Z}
\newcommand{\EXC}{\textup{(EXC)}}

\hypersetup{
  colorlinks=true,
  linkcolor=black,
  citecolor=black,
  urlcolor=blue,
  pdftitle={Strict Cross-Slice Order and Line Batching in the Exact Fisher Network},
  pdfauthor={Maarten Linus Herrmann (ORCID: 0009-0005-8721-6588)},
  pdfsubject={Revised research draft, 27 October 2026},
  pdfkeywords={Fisher exact test; network algorithm; M-concavity; line batching; exact computation},
  pdfcreator={LaTeX}
}

\title{Strict Cross-Slice Order and Line Batching in the Exact Fisher Network}
\author{Maarten Linus Herrmann\\
\small Erasmus School of Economics\\
\small Erasmus University Rotterdam\\
\small\href{https://orcid.org/0009-0005-8721-6588}{ORCID: 0009-0005-8721-6588}}
\date{\small Revised research draft\\
\small 27 October 2026}

\usepackage{xurl}
\begin{document}
\maketitle

\begin{abstract}
At one step of the Mehta--Patel exact Fisher network, a daughter
allocation \(x\) fixes the next table column.  Write
\(B(x)=A(x)L(R-x;C')\) for the largest complete branch weight below it.
Network-induced M-concavity of \(\log L\), strengthened by the strict
first-column multinomial factor, yields a strict cross-slice rhombus.
A within-line ordering lemma then implies monotone movement of leftmost
daughter modes across adjacent three-coordinate slices and an amortized
\(O(m+s)\) peak-query sweep.  The accompanying three-row implementation includes daughter-line
batching, outer-interval acceptance, interval-mass aggregation and prefix
reuse. These are implementation mechanisms distinct from the query bound.
The release keeps exact-rational verification
and FEXACT-compatible floating computation as explicit numerical scopes.
\end{abstract}

\section{Introduction}

The Mehta--Patel network computes the Fisher--Freeman--Halton exact tail
for a contingency table with fixed margins by merging equal residual
states and trimming complete subtrees with exact continuation
bounds~\cite{mehta1983,mehta1986}.  FEXACT adds packed state
representations, shared bound caches and ordered prefix
storage~\cite{clarkson1993}.  Later work develops exact path trimming,
discrete-convex transportation formulations and one-dimensional
daughter ordering~\cite{aoki2002,suzuki2005,requena2006}.

This paper isolates one structural result within a broader implementation.
Sections 2--6 prove a strict cross-slice branch inequality and its
monotone-leftmost-mode consequence.  Section 7 states exactly what that
consequence changes in an implementation.  Section 8 describes the implementation and its finite correctness tests.
The theorem controls peak location; it does not imply reductions in runtime,
arcs, states or memory.

The conditional exact test and network construction are due to Mehta and Patel
\cite{mehta1983,mehta1986}; the practical FEXACT refinements are due to
Clarkson, Fan and Joe~\cite{clarkson1993}; the network-induced
discrete-convex structure used in the proof is due to Murota
\cite{murota1996,murota2021}; and the later path-trimming,
transportation and daughter-ordering developments are due respectively
to Aoki; Suzuki, Aoki and Murota; and Requena and Mart\'in Ciudad
\cite{aoki2002,suzuki2005,requena2006}.  This paper establishes the strict
cross-slice rhombus, its monotone-leftmost-mode and amortized peak-query
consequences, and a research implementation.  The cited literature
supplies the exact test, network algorithm, FEXACT, network induction,
transportation extrema and one-dimensional daughter ordering.

In standard array language, the cross-slice result is a strict multiplicative
Monge-type inequality, equivalently strict log-submodularity, on each
feasible elementary slice rhombus.

\section{Exact Fisher network and branch bounds}

Let \(Y=(y_{ab})\) be a nonnegative \(r\times c\) integer table with
fixed row and column margins.  Conditional probability is proportional
to
\[
  \frac{1}{\prod_{a,b}y_{ab}!}.
\]
The exact two-sided tail used here contains every table no more probable
than the observed table.  At a residual network state, let
\(R\in\Z_+^r\) be the row margins, let the next column total be \(m\),
and let the later column margins be \(C'=(c_1,\ldots,c_q)\).
For a feasible daughter allocation \(x\), define
\[
 A(x)=\binom{m}{x_1,\ldots,x_r},
 \qquad
 L(y;C')=\max_Z\prod_{b=1}^q
 \binom{c_b}{z_{1b},\ldots,z_{rb}},
\]
where \(Z\) has row margins \(y\) and column margins \(C'\).  The exact
largest complete branch weight below \(x\) is
\begin{equation}
  B(x)=A(x)L(R-x;C').
  \label{eq:branch}
\end{equation}
Every value used below is positive and finite because its daughter and
suffix margins are explicitly required to be feasible.

There are two numerical scopes.  The research oracle uses
arbitrary-precision integer path weights, exact tail membership, exact
rational output and replayable certificates.  The native product uses
floating log-factorials, the FEXACT comparison convention with tolerance
\(3.45254\times10^{-7}\), and floating output.  Release tests find the
same conditional tail on their declared cases, but the scopes do not
provide identical numerical guarantees.  Certificate generation and
verification are timed separately from uncertified native computation.

\section{Discrete-convex suffix structure}

The feasible suffix-margin domain is
\[
 \mathcal D(C')=
 \left\{y\in\Z_+^r:\sum_a y_a=\sum_b c_b\right\}.
\]
Because the transportation graph is complete bipartite, this is exactly
the domain of feasible row margins for \(C'\), including vectors with
zero coordinates.

The M-concavity convention in this paper is the following exchange
axiom.  For a finite function \(f\) on a fixed-sum integral domain, if
\(u_a>v_a\), there is a coordinate \(b\) with \(u_b<v_b\) such that
\begin{equation*}
 f(u)+f(v)\le
 f(u-e_a+e_b)+f(v+e_a-e_b).
 \tag{EXC}
\end{equation*}
The exchanged vectors must belong to the domain.  This is equation
(2.4) on printed page 277 of Murota~\cite{murota1996}.

\begin{lemma}[Network-induced suffix M-concavity]
\label{lem:suffix}
The function \(\ell(y)=\log L(y;C')\) is M-concave on
\(\mathcal D(C')\).
\end{lemma}

\begin{proof}
For a suffix table \(Z=(z_{ab})\),
\begin{equation}
 \log L(y;C')=
 \sum_b\log(c_b!)
 -\min_Z\sum_{a,b}\log(z_{ab}!).
 \label{eq:transport}
\end{equation}
Thus \(\ell\) is a constant minus a separable convex transportation
minimum.  More explicitly, map \(z_{ab}\) to integer flow on the
complete bipartite network from row vertices to column vertices.  Its
row boundary is \(y\), its column demand is \(C'\), and each arc has
cost
\[
 \phi(n)=\log(n!),\qquad
 \phi(n+1)-\phi(n)=\log(n+1).
\]
The increments are nondecreasing, so \(\phi\) is discrete convex.
The negative optimum value of this separable convex-cost flow is
M-concave in its full boundary by Murota's Example 2.3~\cite{murota1996};
Theorem 4.7(2) of~\cite{murota2021} also gives the convex form for
separable convex arc costs. Fixing the column boundary preserves the
row exchange: the two column boundary vectors agree, so an exchange
recipient must be a row coordinate.  The first term
of \eqref{eq:transport} is independent of \(y\), proving the claim.
\end{proof}

Zero residual coordinates are permitted by \(\mathcal D(C')\).
An exchange decrements only a strict surplus and cannot create a
negative coordinate.  If \(C'\) is empty, then
\(\mathcal D(C')=\{0\}\), \(L(0;\varnothing)=1\), and the exchange
statement is vacuous.  Neighboring vectors outside the domain are not
assigned approximate or surrogate values; statements below require all
displayed neighbors to be feasible.

\section{Strict within-line order}

The following standard balancing argument makes the mode argument
self-contained.  One-dimensional daughter ordering and its discrete-convex
precedents include Requena and Mart\'in Ciudad~\cite{requena2006}.

\begin{lemma}[Balanced integer split]
\label{lem:split}
Let \(Z\) be a nonnegative integer matrix with even row and column sums,
and suppose specified entries of \(Z\) are even.  Then
\(Z=U+V\) for nonnegative integer matrices \(U,V\) that have half the
row and column sums of \(Z\), satisfy
\(|u_{ab}-v_{ab}|\le1\) in every cell, and split every specified even
entry equally.
\end{lemma}

\begin{proof}
Split every even entry equally.  Make a bipartite graph with one edge
for each odd entry, joining its row and column vertices.  Each vertex
has even degree because its margin and the contributions of even entries
are even.  Decompose the graph into cycles and color each cycle
alternately.  Give the larger half of each odd entry to \(U\) on one
color and to \(V\) on the other.  The two colors occur equally often at
each vertex, which gives the required margins and cellwise balance.
\end{proof}

\begin{lemma}[Strict within-line log-concavity]
\label{lem:line}
Choose distinct rows \(p,q\), set
\(x(t)=x(0)+t(e_p-e_q)\), and let \(I\subset\Z\) be the feasible
integer interval.  For every adjacent triple in \(I\),
\begin{equation}
  B(x(t))^2>B(x(t-1))B(x(t+1)).
  \label{eq:line-strict}
\end{equation}
Consequently the adjacent ratios \(B(x(t+1))/B(x(t))\) strictly
decrease in \(t\).  The line has either one maximum or two adjacent
equal maxima.  Its leftmost maximizer is
\begin{equation}
 p=\min\left(
 \{t\in I:t<\max I,\ B(x(t+1))/B(x(t))\le1\}
 \cup\{\max I\}\right).
 \label{eq:leftmost-characterization}
\end{equation}
\end{lemma}

\begin{proof}
Let \(Y^-\) and \(Y^+\) minimize the products of cell factorials among
complete tables whose first columns are \(x(t-1)\) and \(x(t+1)\).
The matrix \(Z=Y^-+Y^+\) has even margins and first column \(2x(t)\).
Apply Lemma~\ref{lem:split}, specifying the first column, to obtain
complete tables \(U,V\) with first column \(x(t)\).

For two nonnegative integers of fixed sum, their factorial product is
minimized when they are as equal as possible: if \(a\ge b+2\), moving
one unit from \(a\) to \(b\) multiplies \(a!b!\) by
\((b+1)/a<1\).  Applying this cellwise gives
\[
 D(U)D(V)\le D(Y^-)D(Y^+),
 \qquad D(W)=\prod_{a,b}w_{ab}!.
\]
The inequality is strict in first-column row \(p\), whose outer values
are \(x_p(t)-1\) and \(x_p(t)+1\), while both middle values are
\(x_p(t)\).  Hence, writing \(d(x)\) for the minimum of \(D\) at first
column \(x\),
\[
 d(x(t))^2\le D(U)D(V)
 <d(x(t-1))d(x(t+1)).
\]
Taking reciprocals and multiplying by the common squared product of
column factorials proves \eqref{eq:line-strict}.

Dividing \eqref{eq:line-strict} by adjacent positive terms shows that
the forward ratios strictly decrease.  They can therefore cross one
at most once.  A ratio equal to one produces exactly two adjacent equal
maxima; otherwise the maximum is unique.  Before the first ratio at
most one, values strictly increase, and at and after it they do not
increase.  This proves \eqref{eq:leftmost-characterization}, including
intervals of length one.
\end{proof}

\section{Strict cross-slice rhombus}

Choose distinct rows \(i,j,k\).  Suppose all four daughters
\[
 x,\quad x^i=x+e_i-e_j,\quad x^k=x+e_k-e_j,\quad
 x^{ik}=x+e_i+e_k-2e_j
\]
are feasible.  In particular, feasibility of \(x^{ik}\) requires
\(x_j\ge2\).

\begin{theorem}[Strict cross-slice rhombus]
\label{thm:rhombus}
\begin{equation}
 B(x^i)B(x^k)>B(x)B(x^{ik}).
 \label{eq:strict-rhombus}
\end{equation}
\end{theorem}

\begin{proof}
Set \(y=R-x\) and \(y^{ik}=R-x^{ik}\).  Then
\begin{equation}
 y-y^{ik}=e_i+e_k-2e_j.
 \label{eq:forced-vector}
\end{equation}
Relative to \(y^{ik}\), the vector \(y\) has surplus coordinates \(i\)
and \(k\), and \(j\) is the only coordinate with the required opposite
sign.  After selecting surplus coordinate \(i\) in \EXC, the
exchange recipient is therefore forced to be \(j\).

The four suffix margins, written explicitly, are
\[
\begin{array}{c|c}
\text{daughter}&\text{suffix margin}\\ \hline
x&y\\
x^i&y-e_i+e_j\\
x^k&y-e_k+e_j\\
x^{ik}&y-e_i-e_k+2e_j.
\end{array}
\]
Lemma~\ref{lem:suffix} instantiates \EXC\ as
\[
\begin{split}
\ell(y)+\ell(y-e_i-e_k+2e_j)
\le{}&\ell(y-e_i+e_j)\\
&+\ell(y-e_k+e_j).
\end{split}
\]
Exponentiating gives the known, non-strict suffix inequality
\begin{equation}
 L(y-e_i+e_j;C')L(y-e_k+e_j;C')
 \ge L(y;C')L(y-e_i-e_k+2e_j;C').
 \label{eq:weak-suffix}
\end{equation}

Direct cancellation of the four first-column multinomials gives
\begin{equation}
 \frac{A(x^i)A(x^k)}{A(x)A(x^{ik})}
 =\frac{x_j}{x_j-1}.
 \label{eq:strict-factor}
\end{equation}
Feasibility of \(x^{ik}\) implies \(x_j\ge2\), so the factor is
strictly greater than one.  Multiplying \eqref{eq:weak-suffix} and
\eqref{eq:strict-factor} proves \eqref{eq:strict-rhombus}.
\end{proof}

The conceptual separation is exact: known network-induced M-concavity
gives the weak inequality; the first-column multinomial gives the strict
factor; their product gives the strict branch rhombus.

\section{Monotone leftmost modes}

Fix all daughter coordinates outside rows \(i,j,k\).  Let \(d\) be the
allocation total remaining in these three coordinates and write
\[
 t=x_i,\qquad z=x_k,\qquad x_j=d-z-t.
\]
Let \(r_i,r_j,r_k\) be the three residual row capacities.  For an
integer \(z\) with \(0\le z\le r_k\), define
\begin{equation}
 I_z=[\ell_z,u_z]\cap\Z,\qquad
 \ell_z=\max(0,d-z-r_j),\qquad
 u_z=\min(r_i,d-z).
 \label{eq:slice-interval}
\end{equation}
The slice is feasible exactly when \(I_z\ne\varnothing\).  For adjacent
feasible slices \(z,z+1\),
\begin{equation}
 \ell_{z+1}=\max(0,d-z-1-r_j)\le\ell_z,\qquad
 u_{z+1}=\min(r_i,d-z-1)\le u_z.
 \label{eq:endpoints}
\end{equation}

\begin{lemma}[Monotone leftmost mode]
\label{lem:generic-mode}
Let \(b(t,z)>0\) be defined on feasible integer intervals
\(I_z=[\ell_z,u_z]\cap\Z\).  Assume that:
\begin{enumerate}
\item \(b(t,z)^2>b(t-1,z)b(t+1,z)\) for every feasible adjacent triple;
\item the forward ratio
\(\rho_z(t)=b(t+1,z)/b(t,z)\) satisfies
\(\rho_{z+1}(t)<\rho_z(t)\) wherever both ratios exist;
\item \(\ell_{z+1}\le\ell_z\) and \(u_{z+1}\le u_z\).
\end{enumerate}
Then the leftmost maximizer \(p(z)\) weakly decreases:
\[
 p(z+1)\le p(z).
\]
\end{lemma}

\begin{proof}
Strict log-concavity makes \(\rho_z(t)\) strictly decreasing in \(t\).
Thus Lemma~\ref{lem:line}'s characterization
\eqref{eq:leftmost-characterization} applies to each slice, including a
slice of length one and the case of two adjacent equal maxima.

Write \(p=p(z)\) and \(I_{z+1}=[\ell',u']\cap\Z\).
If \(p>u'\), then \(p(z+1)\le u'<p\); this is the upper-endpoint clipping
case.  Otherwise, \(\ell'\le\ell_z\le p\le u'\), so the previous mode
remains feasible.  If \(p=u'\), the desired inequality is immediate.
If \(p<u'\), the ratio \(\rho_{z+1}(p)\) exists.  Because
\(u'\le u_z\), the ratio \(\rho_z(p)\) also exists.  The leftmost-mode
characterization gives \(\rho_z(p)\le1\), with equality precisely when
\(p\) is the left member of two adjacent equal maxima.  Assumption 2
then gives
\[
  \rho_{z+1}(p)<\rho_z(p)\le1.
\]
Therefore the first ratio at most one on slice \(z+1\), or its upper
endpoint if no such ratio exists, occurs no later than \(p\).
Hence \(p(z+1)\le p\).  This also covers ratios at the interval
endpoints: a forward ratio is used only when its right neighbor exists.
\end{proof}

\begin{corollary}[Fisher-network slice order]
\label{cor:fisher-mode}
For \(b(t,z)=B(x)\) under \eqref{eq:slice-interval}, the leftmost
maximizer weakly decreases with \(z\).
\end{corollary}

\begin{proof}
Lemma~\ref{lem:line} gives assumption 1 of
Lemma~\ref{lem:generic-mode}, and \eqref{eq:endpoints} gives assumption
3.  Theorem~\ref{thm:rhombus}, with increments in \(t\) and \(z\), is
\[
 B(t+1,z)B(t,z+1)>B(t,z)B(t+1,z+1).
\]
Dividing by positive terms yields
\[
 \rho_{z+1}(t)=\frac{B(t+1,z+1)}{B(t,z+1)}
 <\frac{B(t+1,z)}{B(t,z)}=\rho_z(t),
\]
which is assumption 2.  Apply Lemma~\ref{lem:generic-mode}.
\end{proof}

\begin{lemma}[Amortized peak-query bound]
\label{lem:query-bound}
Consider \(s\ge 1\) consecutive feasible Fisher-network slices as
in \eqref{eq:slice-interval}, all contained in the common coordinate
range \(\{0,\ldots,m\}\). Find the first leftmost mode by binary search.  On every later slice,
start at the preceding mode clipped to the new upper endpoint and walk
left to the new leftmost mode.  With branch values cached within the
sweep, peak location uses
\[
 O(\log m)+O(m)+O(s)=O(m+s)
\]
branch-bound queries and peak comparisons.  Independent binary searches
use \(O(s\log m)\).
\end{lemma}

\begin{proof}
Write \(p_h\) for the leftmost mode of slice \(h\), and for \(h\ge1\)
let \(q_h=\min(p_{h-1},u_h)\) be the clipped starting point.  The
corollary and feasibility give
\[
 p_h\le q_h\le p_{h-1}.
\]
The number of left moves on slice \(h\) is \(q_h-p_h\), and therefore
\[
 \sum_{h=1}^{s-1}(q_h-p_h)
 \le \sum_{h=1}^{s-1}(p_{h-1}-p_h)
 =p_0-p_{s-1}\le m.
\]
There is one terminal check per later slice, and each ratio comparison
requests only a constant number of cached branch values.  The first
binary search costs \(O(\log m)\).  This proves the displayed bound,
with the cases \(m\le1\) absorbed into the constant.
\end{proof}

Lemma~\ref{lem:query-bound} covers branch-bound calls and peak-location
comparisons.  Daughter generation, threshold searches, accepted intervals,
forwarded arcs, residual states and total runtime depend on separate
mechanisms.  Peak comparisons can therefore fall without a corresponding
wall-clock improvement.

\section{Algorithmic consequence}

The shared native C core exposes five explicit modes:
\texttt{baseline}, \texttt{accept\_only}, \texttt{line\_batch},
\texttt{line\_batch\_monge} and \texttt{auto}.  The two line modes have
identical daughter generation, exact outer-interval membership,
interval-mass aggregation, reduced forwarding and prefix-state reuse.
They differ only in peak location: \texttt{line\_batch} searches each
slice independently, whereas \texttt{line\_batch\_monge} uses
Corollary~\ref{cor:fisher-mode}. The two modes isolate the peak-location change.

The native engine supports \(3\times c\) tables and, by safe
transposition, \(r\times3\) tables.  Other dimensions are explicitly
delegated by the R interface to
\texttt{stats::fisher.test(hybrid=FALSE, simulate.p.value=FALSE)};
the standalone program reports them as unsupported.  The product is not
claimed as a general \(r\times c\) replacement.

In exact arithmetic, accept-only pruning preserves the tail. It omits
shortest-path rejection and retains longest-bound acceptance; every
unresolved daughter is forwarded.  It can expand excessively. The native floating implementation remains
subject to the numerical convention stated above.  The current \texttt{auto} rule selects
\texttt{line\_batch}.

\section{Implementation and verification}

The line-batching implementation combines outer-interval acceptance,
interval-mass aggregation and prefix reuse. These mechanisms are separate
from the monotone-mode sweep. The query bound does not establish a runtime
or memory advantage, and the implementation is not a general replacement
for FEXACT.

The recorded release gate compares the exact-rational oracle, all manual
native modes and \texttt{stats::fisher.test} on 1,198 cases: exhaustive
positive-margin \(3\times2\) tables with total at most 7, \(3\times3\)
tables with total at most 5, and 60 seeded random cases. It also compares
48 representative cases with the pinned R 4.6.1 FEXACT harness, using
absolute and relative tolerances of \(2\times10^{-12}\). Complementary
tests cover permutations, transposition, workspace and allocation failures,
and interruption. These finite comparisons do not prove numerical agreement
for all inputs.

\section{Empirical boundary and failed conjectures}

Fixed-threshold boundaries need not move monotonically and need not move
one unit at a time; retained exact counterexamples begin at total eight.
A separate shared-longest-DP conjecture also failed: a root call already
closes the reachable suffix DAG, so deleting later top-level calls need
not delete internal dynamic-programming states.

The package remains experimental and is limited to three-row native execution.  Integrating the cross-slice
walk into R's FEXACT core would require a larger architectural change:
FEXACT canonicalizes residual rows and emits daughters individually, whereas
the walk needs stable row labels and ordered adjacent slices.

\section{Attribution and scope of contribution}

This work builds on the cited exact-testing, network-algorithm and
discrete-convexity literature.  In particular, the Mehta--Patel network,
Algorithm 643/FEXACT and its refinements, network-induced M-concavity,
transportation formulations, path trimming and established
one-dimensional daughter ordering remain the contributions of their
respective authors~\cite{mehta1983,mehta1986,clarkson1993,murota1996,
murota2021,aoki2002,suzuki2005,requena2006}.  This paper contributes the
strict cross-slice statement and its bounded algorithmic consequence,
together with a research implementation and finite correctness checks.  Citations identify
intellectual provenance; they do not imply endorsement by the cited
authors or institutions.

\section{Reproducibility and computational assistance}

The public implementation corresponding to this paper is available at
\begin{quote}
\scriptsize
\url{https://github.com/Maar10Herr/Exact-Fisher-Network-Public}.
\end{quote}
It contains the shared C core, R and command-line interfaces, the
exact-rational Python reference, and deterministic verification tests.
The public checks are run with \texttt{make test}; the R package is built
and checked separately with \texttt{make check-r}.

The project record reports that AI tools assisted with literature organization,
conjecture exploration, proof drafting, implementation, testing and record
assembly; the earlier model identifiers are not recorded in the reviewed
project material. GPT-6.1 Terra assisted the present manuscript revision, and
GPT-6.1 Sol assisted preparation of code and verification materials. AI output
is not mathematical evidence: the arguments depend on the cited
network-induction result and explicit factorial identities, while the
computational evidence comprises the stated finite tests.

\section{Conclusion}

Known suffix M-concavity and a strict multinomial factor give a strict
cross-slice rhombus for exact Fisher-network branch values. With a
within-line lemma and moving-domain argument, the rhombus orders leftmost
modes across slices and yields an amortized peak-query bound. The theorem
does not establish a runtime or memory improvement.

\begingroup
\footnotesize
\bibliographystyle{plain}
\bibliography{strict_cross_slice_fisher_network}
\endgroup

\end{document}
